% !TeX encoding = UTF-8
% Français de France. Charger un seul archétype depuis le document principal.
\newcommand{\cvsummary}{Ingénieur en mécanique formé à NYU Abu Dhabi et en mécatronique dans le programme EU4M, avec une expérience en sciences spatiales. Association du traitement de données de missions martiennes assisté par intelligence artificielle et apprentissage automatique à une pratique de la commande embarquée, de la conception mécanique et des systèmes robotiques pour intégrer la mécanique à l'Industrie 4.0.}

% --- COMPÉTENCES ---
\newcommand{\skills}{
    \section*{COMPÉTENCES TECHNIQUES}
    {\small
    \textbf{Programmation et IA :} Python (NumPy, SciPy), MATLAB, C++, Java, SQL, HTML/CSS/JS, Git, ROS 2, notions de PyTorch/fastai, traitement d'images, React.

    \vspace{1pt}
    \textbf{Conception mécanique et simulation :} SolidWorks (CSWA), CATIA V5, ANSYS (Mechanical, Fluent, AQWA), COMSOL, analyse par éléments finis (FEA), mécanique des fluides numérique (CFD), analyse vibratoire, cotation dimensionnelle et géométrique (GD\&T), Solid Edge, Moldflow.

    \vspace{1pt}
    \textbf{Robotique et systèmes embarqués :} Microcontrôleurs, FPGA, Raspberry Pi, automates programmables, conception de circuits imprimés (PCB), capteurs, instrumentation, commande PID, Simulink, représentation d'état et LQR, fusion de données inertielles (IMU), machines à états finis.

    \vspace{1pt}
    \textbf{Sciences spatiales et calcul scientifique :} Mécanique orbitale, traitement d'images satellitaires, cartographie planétaire, SIG (PostgreSQL/PostGIS), simulation de vol, intégration numérique, numérisation 3D, visualisation scientifique.

    \vspace{1pt}
    \textbf{Fabrication et matériaux :} Usinage CNC, fabrication additive (PLA/TPU), moulage par injection, composites en fibre de carbone, sélection de matériaux, fabrication MEMS (photolithographie, PVD/pulvérisation cathodique).

    \vspace{1pt}
    \textbf{Langues :} Espagnol langue maternelle, anglais C2, français C2, italien A1, allemand A1.
    }
}


% --- PROJETS ---
\newcommand{\vawttitle}{VAWT -- Récupération d'énergie éolienne urbaine}
\newcommand{\vawtdates}{En cours (date de début non précisée)}
\newcommand{\vawtlocation}{Abou Dabi, ÉAU}
\newcommand{\vawtdescOne}{Développement d'un concept d'éolienne à axe vertical exploitant le vent généré par le trafic, en transformant une source d'écoulement non conventionnelle en conception de récupération d'énergie urbaine.}
\newcommand{\vawtdescTwo}{Évaluation du concept d'éolienne par simulations CFD sous ANSYS Fluent, à l'aide d'analyses numériques des écoulements pour étudier le comportement aérodynamique.}

\newcommand{\fsaetitle}{Formula SAE -- Suspension et châssis en fibre de carbone}
\newcommand{\fsaedates}{Automne 2024 -- Aujourd'hui}
\newcommand{\fsaelocation}{Abou Dabi, ÉAU}
\newcommand{\fsaedescOne}{Conception de la suspension et du châssis en fibre de carbone d'un véhicule Formula SAE, en combinant géométrie sous SolidWorks et analyse par éléments finis sous ANSYS.}
\newcommand{\fsaedescTwo}{Étude de la cinématique de suspension sous SolidWorks et ANSYS, en exploitant le comportement simulé du mécanisme pour orienter le développement des performances dynamiques.}

\newcommand{\tactitle}{Mini-réfrigérateur -- Système thermoélectrique TEC à deux étages}
\newcommand{\tacdates}{Printemps 2025}
\newcommand{\taclocation}{Abou Dabi, ÉAU}
\newcommand{\tacdescOne}{Conception d'un système de refroidissement sans fluide frigorigène autour d'un module TEC à deux étages, en appliquant la thermodynamique à une alternative à l'état solide à la réfrigération classique.}
\newcommand{\tacdescTwo}{Instrumentation de l'ensemble de refroidissement pour un suivi thermique en temps réel, créant une base expérimentale pour observer et optimiser son comportement.}

\newcommand{\orbitaltitle}{OrbitalGrid -- Système SIG pour applications spatiales}
\newcommand{\orbitaldates}{06/2024 - Aujourd'hui}
\newcommand{\orbitallocation}{Abou Dabi, ÉAU}
\newcommand{\orbitaldescOne}{Développement d'un système SIG fondé sur SQL combinant numérisation 3D et données satellitaires pour la cartographie aérospatiale, en reliant géométrie locale et informations issues de l'espace.}
\newcommand{\orbitaldescTwo}{Mise en œuvre de l'ingénierie des données spatiales sous PostgreSQL pour organiser les informations du système cartographique OrbitalGrid.}

\newcommand{\softrobtitle}{Exosquelette robotique pour la rééducation après un AVC}
\newcommand{\softrobdates}{04/2024 - 05/2024}
\newcommand{\softroblocation}{Abou Dabi, ÉAU}
\newcommand{\softrobdescOne}{Conception d'un exosquelette de doigts inspiré de l'origami pour la rééducation après un AVC, avec actionnement par fils SMA pour traduire un concept biomécanique en mécanisme robotique.}
\newcommand{\softrobdescTwo}{Association de polymère Dragon Skin et de fils en alliage à mémoire de forme pour développer une structure d'actionnement souple destinée à reproduire la fonction musculaire des doigts.}

\newcommand{\exoplanettitle}{Simulateur de transits d'exoplanètes}
\newcommand{\exoplanetdates}{Printemps 2025}
\newcommand{\exoplanetlocation}{Abou Dabi, ÉAU}
\newcommand{\exoplanetdescOne}{Développement en Python d'un simulateur de transits d'exoplanètes générant des séries temporelles de courbes de lumière à partir de paramètres de mécanique orbitale et de physique stellaire.}
\newcommand{\exoplanetdescTwo}{Exploration de différentes configurations étoile--planète par paramétrage des données physiques d'entrée du modèle de transit.}

\newcommand{\windTurbinetitle}{Éolienne flottante en mer -- Stabilisation active}
\newcommand{\windTurbinedates}{05/2026 - 06/2026}
\newcommand{\windTurbinelocation}{Besançon, France}
\newcommand{\windTurbinedescOne}{Conception mécanique d'une plateforme offshore semi-submersible à trois colonnes avec un système de commande pour stabiliser un prototype d'éolienne à échelle réduite, en analysant le tangage et le roulis au moyen d'un MPU6050 et d'une commande en boucle fermée par servomoteurs pour la stabilisation active de la nacelle.}

\newcommand{\carbonmetertitle}{CARBONMETRE -- Écoconception et mesure carbone}
\newcommand{\carbonmeterdates}{09/2025 - 12/2025}
\newcommand{\carbonmeterlocation}{Besançon, France}
\newcommand{\carbonmeterdescOne}{Association de la fabrication, de l'écoconception et de la métrologie dans le projet CARBONMETRE, en appliquant les principes de l'ingénierie de la mesure à l'évaluation carbone.}
\newcommand{\carbonmeterdescTwo}{Application des principes de fabrication industrielle dans un projet de mesure carbone, en reliant les travaux de fabrication à l'ingénierie environnementale.}

\newcommand{\hvactitle}{Amélioration thermique d'une chambre en résidence -- Refroidissement passif et évaporatif}
\newcommand{\hvacdates}{06/2026}
\newcommand{\hvaclocation}{Besançon, France}
\newcommand{\hvacdescOne}{Réduction déclarée de la température d'une chambre de 9 m$^2$ jusqu'à un maximum de 3 \textdegree{}C par rapport à la situation avant intervention, grâce à un écran solaire extérieur et à un refroidisseur à eau gelée entraîné par ventilateur.}
\newcommand{\hvacdescTwo}{Blocage du rayonnement solaire direct pendant la période de pointe de 14 h à 21 h grâce à un écran extérieur de papier aluminium réfléchissant, de PSE de 1 cm et de carton ondulé.}

\newcommand{\controlBoxtitle}{Installation électrique complète -- Équipements pour événements de taille moyenne}
\newcommand{\controlBoxdates}{06/2023 - 08/2023}
\newcommand{\controlBoxlocation}{Mexico, Mexique}
\newcommand{\controlBoxdescOne}{Réduction de 30\% du temps de montage du système électrique par rapport à l'installation précédente, grâce à la conception d'un coffret de commande de 12 entrées pour des événements de taille moyenne.}
\newcommand{\controlBoxdescTwo}{Conception d'un système électrique à automate programmable pour des événements en direct, en traduisant les exigences de fonctionnement en matériels de commande coordonnés.}

\newcommand{\hyperlooptitle}{NYC Hyperloop -- Prototype à échelle réduite}
\newcommand{\hyperloopdates}{01/2024 - 05/2024}
\newcommand{\hyperlooplocation}{New York, États-Unis}
\newcommand{\hyperloopdescOne}{Pilotage de l'équipe mécanique d'un Hyperloop à échelle réduite pour la liaison New York--Long Island, avec suspension électrodynamique Inductrack et coordination dans le cadre de la conception globale.}
\newcommand{\hyperloopdescTwo}{Étude de la lévitation magnétique sous COMSOL, ANSYS et SolidWorks par simulation du principe de fonctionnement du système de transport.}

\newcommand{\dronetitle}{Quadricoptère fonctionnel -- Structure allégée}
\newcommand{\dronedates}{03/2024 - 05/2024}
\newcommand{\dronelocation}{New York, États-Unis}
\newcommand{\dronedescOne}{Réduction de 40\% du poids du quadricoptère par rapport à la référence du projet grâce à une structure imprimée en 3D en PLA et TPU.}
\newcommand{\dronedescTwo}{Évaluation du comportement aux chocs par essais de chute selon les critères de performance du projet, en associant essais physiques et développement de la structure.}

\newcommand{\trentxwbtitle}{Rolls-Royce Trent XWB -- Modèle 3D fonctionnel}
\newcommand{\trentxwbdates}{01/2024 - 05/2024}
\newcommand{\trentxwblocation}{New York, États-Unis}
\newcommand{\trentxwbdescOne}{Codéveloppement d'un modèle 3D complet du Rolls-Royce Trent XWB sous SolidWorks, en transformant un système de propulsion complexe en assemblage CAO, avec analyse structurelle de l'arbre principal sous SolidWorks et simulations aérodynamiques sous ANSYS.}
\newcommand{\trentxwbdescTwo}{Étude du comportement aérodynamique sous ANSYS, en prolongeant le projet de modélisation du moteur au-delà de la représentation géométrique.}

\newcommand{\liostitle}{LIOS -- Assistant vocal avec IA}
\newcommand{\liosdates}{12/2024 - Aujourd'hui}
\newcommand{\lioslocation}{Abou Dabi, ÉAU}
\newcommand{\liosdescOne}{Développement d'un assistant vocal Python combinant traitement automatique du langage, données en temps réel et API externes, en transformant des fonctions d'IA appliquée en système logiciel interactif.}
\newcommand{\liosdescTwo}{Mise en place d'une mémoire contextuelle conservant les informations au fil des interactions de l'assistant, pour étendre les fonctions au-delà des requêtes isolées.}

\newcommand{\fitpaltitle}{FitPal -- Dispositif intelligent d'entraînement physique}
\newcommand{\fitpaldates}{09/2021 - 12/2021}
\newcommand{\fitpallocation}{Abou Dabi, ÉAU}
\newcommand{\fitpaldescOne}{Codéveloppement d'un dispositif intelligent d'entraînement physique sur M5Stack pour l'analyse d'exercices en temps réel, intégrant matériel embarqué et mesures destinées à l'utilisateur.}
\newcommand{\fitpaldescTwo}{Mise en place du transfert de données IoT et de la connectivité mobile via Arduino IoT App, en reliant le dispositif à une interface distante.}

\newcommand{\covidtitle}{Cabine de désinfection automatisée}
\newcommand{\coviddates}{09/2020 - 02/2021}
\newcommand{\covidlocation}{Mexico, Mexique}
\newcommand{\coviddescOne}{Conception d'un prototype de cabine de désinfection automatisée sous Arduino, intégrant électronique de commande et système physique de brumisation.}
\newcommand{\coviddescTwo}{Mise en place d'une communication sans fil dans le prototype de cabine, en ajoutant une connectivité distante au système de commande.}

\newcommand{\rovertitle}{Rover téléopéré -- Numérisation spatiale 3D}
\newcommand{\roverdates}{09/2026 - Aujourd'hui}
\newcommand{\roverlocation}{Oviedo, Espagne}
\newcommand{\roverdescOne}{Développement en cours d'un rover téléopéré pour la numérisation spatiale 3D, intégrant commande embarquée, acquisition de nuages de points par capteur flexx2 et visualisation distante via communication TCP/IP entre M5Stack et Raspberry Pi ; conception du châssis et de la structure mécanique pour loger l'électronique dans un ensemble en ABS moulé par injection.}
\newcommand{\roverdescTwo}{Développement en cours d'un rover de numérisation 3D combinant commande embarquée, acquisition de nuages de points et visualisation/commande distantes ; conception du châssis et de la structure mécanique pour moulage par injection d'ABS.}

\newcommand{\microsensortitle}{Microcapteur piézorésistif de pression de détonation}
\newcommand{\microsensordates}{05/2026 - 06/2026}
\newcommand{\microsensorlocation}{Besançon, France}
\newcommand{\microsensordescOne}{Réduction des itérations de géométrie et de dessin des masques par automatisation en Python du dimensionnement de capteurs piézorésistifs, en remplaçant les recalculs manuels par des contraintes de K, $\Delta R/R_0$ et résistivité imposées par le code.}
\newcommand{\microsensordescTwo}{Pilotage de la sélection de matériaux d'un capteur à l'échelle du wafer visant une interface de 50 $\Omega$ et une tenue à 10 GPa, en conciliant exigences électriques, mécaniques et économiques couplées.}

\newcommand{\missiletitle}{Simulateur de dynamique d'interception de missiles}
\newcommand{\missilelocation}{Besançon, France}
\newcommand{\missiledescOne}{Développement sous MATLAB d'une simulation 3D d'interception de missiles avec loi de guidage par navigation proportionnelle, intégration numérique et visualisation des trajectoires en temps réel avec détection d'impact.}

% --- EXPÉRIENCE ---
\newcommand{\carbonorg}{Energy Innovation Capital}
\newcommand{\carbontitle}{Analyste de marché -- Captage du carbone}
\newcommand{\carbondates}{10/2023 - 01/2024}
\newcommand{\carbonlocation}{Abou Dabi, ÉAU}
\newcommand{\carbonbullets}{%
  \item Élaboration de recommandations d'investissement CCUS pour un fonds de capital-risque d'Abou Dabi, en transformant les recherches sur le captage du carbone en évaluations de jeunes entreprises.
  \item Synthèse des avancées scientifiques et techniques du captage du carbone en présentations pour les analystes seniors, en reliant les résultats de recherche aux questions d'investissement.
  \item Production d'analyses des tendances du marché reliant moteurs technologiques et risques d'investissement, par examen des données du secteur du captage du carbone.
  \item Mise en relation de la recherche technique et de l'évaluation commerciale par recommandations stratégiques fondées sur la science CCUS et l'analyse de marché.
}

\newcommand{\spacelaborg}{NYU Abu Dhabi, Centre des sciences spatiales}
\newcommand{\spacelabtitle}{Ingénieur et chercheur R\&D -- Plateforme de sciences planétaires}
\newcommand{\spacelabdates}{05/2025 - 09/2025}
\newcommand{\spacelablocation}{Abou Dabi, ÉAU}
\newcommand{\spacelabbullets}{%
  \item Déploiement d'un atlas interactif de Mars fondé sur des images multispectrales réelles d'Emirates Mars Mission HOPE, avec développement de l'\href{https://romvio.github.io/mars-viewer/}{application web et mobile} en React et JavaScript et soumission à l'App Store d'Apple.
  \item Amélioration de 40\% du traitement des images planétaires par compression optimisée et correction des pixels assistée par apprentissage automatique, en assemblant les images HOPE par calibration des recouvrements, mosaïquage et normalisation de la luminosité.
  \item Collecte d'échantillons de sols et de sédiments sur des sites désertiques des ÉAU par carottage au tube et excavation, avec documentation du terrain et des sédiments pour rechercher de possibles analogues du régolithe martien.
}

\newcommand{\geoorg}{Global Education NYUAD}
\newcommand{\geotitle}{Chargé de coordination Global Education -- Étudiants en mobilité entrante}
\newcommand{\geodates}{12/2025 - 05/2025}
\newcommand{\geolocation}{Abou Dabi, ÉAU}
\newcommand{\geobullets}{%
  \item Coordination de sorties universitaires et culturelles aux ÉAU pour les étudiants internationaux en mobilité entrante, en transformant les plans du programme en activités de groupe organisées.
  \item Pilotage de la planification logistique de l'équipe de coordinateurs et de la répartition des responsabilités pour accompagner des groupes étudiants variés.
  \item Intégration de protocoles de sécurité à la préparation des sorties, en reliant la logistique de groupe aux précautions opérationnelles.
  \item Rôle de référent des étudiants en mobilité, en les mettant directement en relation avec les services d'accompagnement universitaires.
}

\newcommand{\geoOneorg}{Global Education NYUAD}
\newcommand{\geoOnetitle}{Chargé de coordination Global Education -- Étudiants de première année}
\newcommand{\geoOnedates}{09/2026 - 12/2024}
\newcommand{\geoOnelocation}{Abou Dabi, ÉAU}
\newcommand{\geoOnebullets}{%
  \item Animation de séances d'accueil de première année, en transformant les informations universitaires en activités d'intégration structurées.
  \item Mentorat et conseil continus entre pairs pour faciliter la transition vers l'université.
  \item Conduite de visites du campus présentant les installations et ressources aux nouveaux étudiants avec des explications directes et adaptées.
  \item Planification d'événements universitaires et culturels inclusifs, créant des occasions d'intégration à la communauté pour les étudiants de première année.
}

\newcommand{\hcorg}{Centre de santé de NYUAD}
\newcommand{\hctitle}{Emploi étudiant -- Responsable des réseaux sociaux}
\newcommand{\hcdates}{09/2024 - 05/2025}
\newcommand{\hclocation}{Abou Dabi, ÉAU}
\newcommand{\hcbullets}{%
  \item Réalisation de campagnes de santé ciblées pour les étudiants, en transformant les initiatives de bien-être en messages adaptés à la communauté de NYU Abu Dhabi.
  \item Conception et publication de contenus pour les réseaux sociaux promouvant les ressources de santé, en transformant les informations de service en supports numériques accessibles.
  \item Rôle de contact principal pour les demandes étudiantes, en les orientant vers les services d'accompagnement du centre de santé.
  \item Contribution à la résolution des demandes par traitement direct, en associant communication publique et aide individuelle aux étudiants.
}

\newcommand{\alpesorg}{Alpes Technologies México}
\newcommand{\alpestitle}{Ingénieur projets junior -- Gestion d'entrepôt et systèmes électriques}
\newcommand{\alpesdates}{06/2024 - 09/2024}
\newcommand{\alpeslocation}{Mexico, Mexique}
\newcommand{\alpesbullets}{%
  \item Amélioration de 60\% de l'efficacité opérationnelle et de fabrication par rapport à un inventaire entièrement manuel, grâce à un système automatique de gestion d'entrepôt sur mesure avec codes QR en HTML, JS et Google Drive.
  \item Mise en place d'une traçabilité à la pièce par mises à jour bidirectionnelles des stocks en temps réel, droits d'accès par rôle et rapports automatisés dans le WMS.
  \item Amélioration du facteur de puissance industriel jusqu'à 99\% par fabrication et installation de batteries de condensateurs pour la compensation harmonique, réduisant les coûts énergétiques des clients de 15\% et les arrêts électriques de 25\% grâce à l'analyse des postes.
  %\item Production de schémas unifilaires et de recommandations correctives pour un complexe hôtelier sinistré par un incendie, par mesures sur site et présentation des résultats aux directions de l'entreprise et du client.
}

\newcommand{\epigsaorg}{EPIGSA}
\newcommand{\epigsatitle}{Assistant en ingénierie des données -- SIG et hydraulique}
\newcommand{\epigsadates}{05/2024 - 07/2024}
\newcommand{\epigsalocation}{Mexico, Mexique}
\newcommand{\epigsabullets}{%
  \item Réduction de 25\% des erreurs de données d'infrastructures par automatisation des contrôles qualité pour des projets municipaux de traitement des eaux.
  \item Intégration et validation de jeux de données SIG sous PostgreSQL, avec préparation d'informations géospatiales pour des infrastructures municipales de traitement des eaux.
  %\item Amélioration de la fiabilité des informations des projets par contrôles automatisés et validation des données SIG intégrées à la base.
  %\item Mise en relation de l'ingénierie des données spatiales et des infrastructures hydrauliques par intégration sous PostgreSQL et contrôle qualité automatisé dans les processus de projet.
}

\newcommand{\solarshiporg}{Solar Ship}
\newcommand{\solarshiptitle}{Stagiaire en stratégie et conception de mobilité aérienne}
\newcommand{\solarshipdates}{05/2023 - 08/2023}
\newcommand{\solarshiplocation}{Toronto, Canada}
\newcommand{\solarshipbullets}{%
  \item Développement de concepts d'infrastructures de logistique aérienne pour des zones isolées de Zambie et de Tanzanie, en reliant conception aérospatiale et besoins d'accès hors réseau.
  \item Contribution à l'initiative Peace+Freedom par un système aéronautique entièrement électrique conçu autour des besoins de connectivité de communautés isolées.
  \item Mise en relation du développement des systèmes d'aéronefs et de la planification des infrastructures, en abordant la logistique aérienne au-delà du seul véhicule.
  %\item Association de l'ingénierie aérospatiale et d'objectifs de durabilité sociale dans un projet multinational pour des communautés hors réseau.
}

\newcommand{\crousorg}{CROUS Besançon Bourgogne Franche-Comté}
\newcommand{\croustitle}{Emploi étudiant -- Restauration universitaire}
\newcommand{\crousdates}{10/2024 - Aujourd'hui}
\newcommand{\crouslocation}{Besançon, France}
\newcommand{\crousbullets}{%
  \item Maintien du rythme en cuisine aux heures de pointe par appui à la préparation alimentaire en grand volume et aux opérations internes.
  \item Respect des pratiques d'hygiène et de sécurité alimentaire par nettoyage, désinfection et plonge en cuisine.
  \item Contribution à la continuité du service par réapprovisionnement pendant les opérations de restauration universitaire.
  \item Coordination des tâches de cuisine et d'appui au service dans un environnement à délais courts, avec adaptation aux besoins opérationnels changeants.
}

\newcommand{\rallyorg}{Championnat du monde des rallyes -- Guanajuato}
\newcommand{\rallytitle}{Commissaire technique officiel FIA}
\newcommand{\rallydates}{11/2019 - 04/2020}
\newcommand{\rallylocation}{Guanajuato, Mexique}
\newcommand{\rallybullets}{%
  \item Marquage officiel et inspection des pneumatiques au Championnat du monde des rallyes FIA à Guanajuato, selon les procédures prescrites de contrôle technique.
  \item Contribution à la conformité des véhicules de compétition par vérifications techniques au regard des exigences de l'épreuve.
  \item Coordination des inspections avec les commissaires responsables et les pilotes sous la pression d'un week-end de compétition.
  \item Application de contrôles qualité aux inspections techniques par procédures de marquage et de vérification décrites dans la base de données conformément aux normes ISO 9001.
}

\newcommand{\papaloteorg}{Papalote Museo del Niño}
\newcommand{\papalotetitle}{Chargé d'accueil et de médiation}
\newcommand{\papalotedates}{05/2019 - 07/2019}
\newcommand{\papalotelocation}{Mexico, Mexique}
\newcommand{\papalotebullets}{%
  \item Présentation d'expositions scientifiques et techniques aux enfants et aux adultes, en transformant les concepts complexes en explications accessibles aux visiteurs.
  \item Proposition d'améliorations concrètes de l'expérience de visite, mises en œuvre pendant l'exploitation du musée.
  \item Adaptation des explications techniques à un public diversifié pour favoriser la participation de personnes d'âges et de besoins différents.
  \item Mise en relation de la médiation scientifique et de l'interaction directe avec le public par découverte guidée des expositions du musée.
}
